\section{Weekly Summary}

\begin{tcolorbox}[colback=yellow!10,colframe=black!50,title=AI-Generated Content]
\noindent The summary below was written by Claude (Anthropic) from that week's regression outputs and series descriptions. It has not been reviewed or fact-checked by a human, and should be read as an automated, informal recap -- not analysis.
\end{tcolorbox}

Welcome back to the weekly ritual where we throw two unrelated economic series into a blender, hit the button, and solemnly report whatever sludge comes out. This week the dartboard gifted us three pairings, and in a shocking twist, the universe decided to be boring, dramatic, and enlightening in exactly that order.

Monday kicked us off with the single most specific series name I have ever encountered: the 90\textbackslash{}\% Confidence Interval Upper Bound of the Estimate of People Under Age 5 in Poverty for the District of Columbia, versus housing units authorized-but-not-started in buildings with two to four units. Just typing that out aged me. And after all that buildup, the regression produced a p-value of 0.88 and an R-squared of basically zero -- the statistical equivalent of a wet handshake. The detrended version was, if anything, even more of a nothingburger. So no, the upper confidence bound of toddler poverty in DC does not predict whether America is dithering on its duplexes. I know, I'm devastated too. Sometimes two random things are just two random things, which, let the record show, is the expected outcome here roughly always.

Tuesday is where it gets juicy, in the sad way this project loves. We paired Exports of Telecommunications and Computer Services against the Import Price Index for Latin American Apparel Manufacturing -- socks and software, truly a dream team. The raw regression came out gloriously "significant," with an R-squared around 0.59 and a p-value with eleven zeros in front of it. For about four seconds you could convince yourself that selling cloud storage abroad somehow moves the price of imported t-shirts. Then we detrended, the R-squared face-planted to 0.006, and the p-value sauntered up to 0.59. Classic. This is the trend mirage in its natural habitat: two series that both happened to drift over the years, holding hands as they wandered upward, mistaken by the raw math for soulmates. Strip the shared drift and the romance evaporates on contact. Chef's kiss.

And then Wednesday, the one that actually behaves. The 1-Month Treasury yield versus the Bank Prime Loan Rate posted an R-squared of 0.98 that -- hold onto something -- survived detrending at basically 0.98, with a slope hovering near one. This one is real, and it's real because these are both short-term interest rates that march to the Fed's drum; of course they move together. So congratulations to this week's one genuinely robust relationship, discovered by randomly pairing two of the most obviously related numbers in all of finance. Riveting detective work on our part. As always, the standing disclaimer applies: these are randomly matched series, I am making zero causal claims, and spurious correlation remains the house specialty. See you next week for more statistically significant nonsense.
